% Part: normal-modal-logic
% Chapter: tableaux
% Section: rules-for-K

\documentclass[../../../include/open-logic-section]{subfiles}

\begin{document}

\olfileid{nml}{tab}{rul}

\olsection{\Ax{K} માટેના નિયમો}

સામાન્ય વિધાનસંબંધી જોડાણીઓ માટેના નિયમો સામાન્ય વિધાનસંબંધી
ચિહ્નિત !!{tableau}s જેવા જ છે; ફક્ત ઉપસર્ગો ઉમેરાય છે.
દરેક કિસ્સામાં, ચિહ્નિત !!{formula}
$\sFmla{S}{!A}[\sigma]$ પર લાગુ થયેલો નિયમ નવાં !!{formula}s
આપે છે, જેમને પણ~$\sigma$ ઉપસર્ગ હોય છે. આ સહજ છે: દાખલા તરીકે,
જો $!A \land !B$ એ~$\sigma$ નામના જગત પર સત્ય હોય, તો
$!A$ અને $!B$ પણ~$\sigma$ પર જ સત્ય છે (અન્ય જગત પર જરૂરી
નથી). વિધાનસંબંધી નિયમો \olref{tab:prop-rules}માં એકત્ર કર્યા છે.

\begin{table}
  \[\def\arraystretch{3}\begin{array}{|c|c|}
    \hline
    \AxiomC{\sFmla{\True}{\lnot !A}[\sigma]}
    \RightLabel{\TRule{\True}{\lnot}}
    \UnaryInfC{\sFmla{\False}{!A}[\sigma]}
    \DisplayProof
    &
    \AxiomC{\sFmla{\False}{\lnot !A}[\sigma]}
    \RightLabel{\TRule{\False}{\lnot}}
    \UnaryInfC{\sFmla{\True}{!A}[\sigma]}
    \DisplayProof
    \\[1ex]
    \hline
    \AxiomC{\sFmla{\True}{!A \land !B}[\sigma]}
    \RightLabel{\TRule{\True}{\land}}
    \UnaryInfC{\sFmla{\True}{!A}[\sigma]}
    \noLine
    \UnaryInfC{\sFmla{\True}{!B}[\sigma]}
    \DisplayProof
    &
    \AxiomC{\sFmla{\False}{!A \land !B}[\sigma]}
    \RightLabel{\TRule{\False}{\land}}
    \UnaryInfC{$\sFmla{\False}{!A}[\sigma] \quad \mid \quad
      \sFmla{\False}{!B}[\sigma]$}
    \DisplayProof
    \\[2ex]
    \hline
    \AxiomC{\sFmla{\True}{!A \lor !B}[\sigma]}
    \RightLabel{\TRule{\True}{\lor}}
    \UnaryInfC{$\sFmla{\True}{!A}[\sigma] \quad \mid \quad
      \sFmla{\True}{!B}[\sigma]$}
    \DisplayProof
    &
    \AxiomC{\sFmla{\False}{!A \lor !B}[\sigma]}
    \RightLabel{\TRule{\False}{\lor}}
    \UnaryInfC{\sFmla{\False}{!A}[\sigma]}
    \noLine
    \UnaryInfC{\sFmla{\False}{!B}[\sigma]}
    \DisplayProof
    \\[2ex]
    \hline
    \AxiomC{\sFmla{\True}{!A \lif !B}[\sigma]}
    \RightLabel{\TRule{\True}{\lif}}
    \UnaryInfC{$\sFmla{\False}{!A}[\sigma] \quad \mid
      \quad \sFmla{\True}{!B}[\sigma]$}
    \DisplayProof
    &
    \AxiomC{\sFmla{\False}{!A \lif !B}[\sigma]}
    \RightLabel{\TRule{\False}{\lif}}
    \UnaryInfC{\sFmla{\True}{!A}[\sigma]}
    \noLine
    \UnaryInfC{\sFmla{\False}{!B}[\sigma]}
    \DisplayProof
    \\[2ex]
    \hline
  \end{array}\]
  \caption{વિધાનસંબંધી જોડાણીઓ માટે ઉપસર્ગવાળા !!{tableau} નિયમો}
  \ollabel{tab:prop-rules}
\end{table}

બંધતાની શરત સામાન્ય !!{tableau}s જેવી જ છે, પરંતુ અહીં માત્ર
!!{formula}s જ નહિ, ઉપસર્ગો પણ સરખા હોવા જોઈએ. તેથી કોઈ શાખામાં
નીચેના બંને હોય તો તે બંધ છે:
\[
\sFmla{\True}{!A}[\sigma] \quad\text{અને}\quad \sFmla{\False}{!A}[\sigma]
\]
કોઈ ઉપસર્ગ $\sigma$ અને !!{formula}~$!A$ માટે.

ધારણાઓ ગોઠવવાના નિયમો પણ સામાન્ય !!{tableau}s જેવા જ છે;
ફક્ત ધારણાઓ માટે આપણે હંમેશા ઉપસર્ગ~$1$ લઈએ છીએ.
(કયો ઉપસર્ગ લઈએ તે મહત્ત્વનું નથી, જો બધી ધારણાઓ માટે તે એક જ હોય.)
દાખલા તરીકે, આપણે કહીએ છીએ કે
\[
!B_1, \dots, !B_n \Proves !A
\]
જો અને માત્ર જો નીચેની ધારણાઓ માટે બંધ ટેબ્લો હોય:
\[
\sFmla{\True}{!B_1}[1], \dots, \sFmla{\True}{!B_n}[1],
\sFmla{\False}{!A}[1].
\]

\iftag{prvBox}{\iftag{prvDiamond}{$\Box$ અને }{$\Box$}}{}%
\iftag{prvDiamond}{$\Diamond$}{} જેવા મોડલ સંક્રિયકો માટે,
ઉપસર્ગ~$\sigma$વાળા !!a{formula} પર લાગુ થયેલા નિયમના
નિષ્કર્ષનો ઉપસર્ગ $\sigma.n$ હોય છે. જોકે કયું $n$ લઈ શકાય
તે ચિહ્ન~$\True$ છે કે~$\False$ તેના પર આધાર રાખે છે.

\iftag{prvBox}{$\TRule{\True}{\Box}$ નિયમ
  $\sFmla{\True}{\Box !A}[\sigma]$ ધરાવતી શાખાને
  $\sFmla{\True}{!A}[\sigma.n]$થી વિસ્તારે છે.
  \iftag{prvDiamond}{તેવી જ રીતે, }{}}{}%
\iftag{prvDiamond}{$\TRule{\False}{\Diamond}$ નિયમ
  $\sFmla{\False}{\Diamond !A}[\sigma]$ ધરાવતી શાખાને
  $\sFmla{\False}{!A}[\sigma.n]$થી વિસ્તારે છે.}{}
\iftag{notprvBox,notprvDiamond}{આ નિયમ}{આ નિયમો} ફક્ત એવા
ઉપસર્ગ~$\sigma.n$ માટે લાગુ કરી શકાય છે જે સંબંધિત શાખા પર
\emph{પહેલેથી} આવે છે. આવા ઉપસર્ગને તે શાખા પર ``વપરાયેલો'' કહીએ.

\iftag{prvBox}{$\TRule{\False}{\Box}$ નિયમ
  $\sFmla{\False}{\Box !A}[\sigma]$ ધરાવતી શાખાને
  $\sFmla{\False}{!A}[\sigma.n]$થી વિસ્તારે છે.
  \iftag{prvDiamond}{તેવી જ રીતે, }{}}{}%
\iftag{prvDiamond}{$\TRule{\True}{\Diamond}$ નિયમ
  $\sFmla{\True}{\Diamond !A}[\sigma]$ ધરાવતી શાખાને
  $\sFmla{\True}{!A}[\sigma.n]$થી વિસ્તારે છે.}{}
જોકે \iftag{notprvBox,notprvDiamond}{આ નિયમ}{આ નિયમો} ફક્ત એવા
ઉપસર્ગ~$\sigma.n$ માટે લાગુ કરી શકાય છે જે સંબંધિત શાખા પર
\emph{પહેલેથી આવતો નથી}. આવા ઉપસર્ગોને તે શાખા માટે ``નવા'' કહીએ.

નિયમો \olref{tab:rules-K}માં આપ્યા છે.

\begin{table}
  \begin{center}
    \def\arraystretch{3}\def\fCenter{}
    \begin{tabular}{|c|c|}
    \hline
    \iftag{prvBox}{
      \AxiomC{\sFmla{\True}{\Box !A}[\sigma]}
      \RightLabel{\TRule{\True}{\Box}}
      \UnaryInfC{\sFmla{\True}{!A}[\sigma.n]}
      \DisplayProof
      &
      \AxiomC{\sFmla{\False}{\Box !A}[\sigma]}
      \RightLabel{\TRule{\False}{\Box}}
      \UnaryInfC{\sFmla{\False}{!A}[\sigma.n]}
      \DisplayProof\\
      $\sigma.n$ વપરાયેલો છે & $\sigma.n$ નવો છે
      \\[1ex]
      \hline}{}
    \iftag{prvDiamond}{
      \AxiomC{\sFmla{\True}{\Diamond !A}[\sigma]}
      \RightLabel{\TRule{\True}{\Diamond}}
      \UnaryInfC{\sFmla{\True}{!A}[\sigma.n]}
      \DisplayProof
      &
      \AxiomC{\sFmla{\False}{\Diamond !A}[\sigma]}
      \RightLabel{\TRule{\False}{\Diamond}}
      \UnaryInfC{\sFmla{\False}{!A}[\sigma.n]}
      \DisplayProof\\
      $\sigma.n$ નવો છે & $\sigma.n$ વપરાયેલો છે
      \\[1ex]
      \hline}{}
    \end{tabular}
  \end{center}
  \caption{\Ax{K} માટેના મોડલ નિયમો.}
  \ollabel{tab:rules-K}
\end{table}

\iftag{prvBox}{\TRule{\True}{\Box}}{\TRule{\False}{\Diamond}}
માટેનો ઉપસર્ગ વપરાયેલો હોવો જોઈએ એવી મર્યાદા જરૂરી છે;
નહિ તો નીચેનું પણ બંધ !!{tableau} ગણવું પડે:
\iftag{notprvBox,notprvDiamond}{%
  \iftag{prvBox}{%
  \begin{oltableau}
    [\pFmla{\True}{\Box \formula{A}}{1}, just = \TAss
      [\pFmla{\False}{\lnot\Box\lnot \formula{A}}{1}, just = \TAss
        [\pFmla{\True}{\formula{A}}{1.1}, just = {\TRule{\True}{\Box}[1]}
          [\pFmla{\True}{\Box\lnot\formula{A}}{1},
            just ={\TRule{\False}{\lnot}[2]}
            [\pFmla{\True}{\lnot\formula{A}}{1.1},
              just = {\TRule{\True}{\Box}[4]}
              [\pFmla{\False}{\formula{A}}{1.1},
                  just ={\TRule{\True}{\lnot}[5]}, close]
            ]
          ]
        ]
      ]
    ]
  \end{oltableau}
  }{
  \begin{oltableau}
    [\pFmla{\True}{\lnot\Diamond\lnot \formula{A}}{1}, just = \TAss
      [\pFmla{\False}{\Diamond\formula{A}}{1}, just = \TAss
        [\pFmla{\False}{\formula{A}}{1.1}, just = {\TRule{\False}{\Diamond}[2]}
          [\pFmla{\False}{\Diamond\lnot\formula{A}}{1},
            just ={\TRule{\True}{\lnot}[1]}
            [\pFmla{\False}{\lnot\formula{A}}{1.1},
              just = {\TRule{\False}{\Diamond}[4]}
              [\pFmla{\True}{\formula{A}}{1.1},
                  just ={\TRule{\True}{\lnot}[5]}, close]
            ]
          ]
        ]
      ]
    ]
  \end{oltableau}
}}{
  \begin{oltableau}
    [\pFmla{\True}{\Box \formula{A}}{1}, just = \TAss
      [\pFmla{\False}{\Diamond \formula{A}}{1}, just = \TAss
        [\pFmla{\True}{\formula{A}}{1.1}, just = {\TRule{\True}{\Box}[1]}
          [\pFmla{\False}{\formula{A}}{1.1},
            just = {\TRule{\False}{\Diamond}[2]}, close]
        ]
      ]
    ]
\end{oltableau}}

પરંતુ $\Box \formula{A} \Entails/ \Diamond \formula{A}$,
તેથી આપણી પ્રમાણપદ્ધતિ અવિશ્વસનીય બની જાય. તેવી જ રીતે,
$\Diamond \formula{A} \Entails/
\Box \formula{A}$ પણ નથી; છતાં
\iftag{prvBox}{\TRule{\False}{\Box}}{\TRule{\True}{\Diamond}}
માટેનો ઉપસર્ગ નવો હોવાની મર્યાદા ન હોય, તો નીચેનું બંધ ટેબ્લો થાય:
\iftag{notprvBox,notprvDiamond}{%
  \iftag{prvBox}{%
  \begin{oltableau}
    [\pFmla{\True}{\lnot\Box\lnot \formula{A}}{1}, just = \TAss
      [\pFmla{\False}{\Box \formula{A}}{1}, just = \TAss
        [\pFmla{\False}{\formula{A}}{1.1}, just = {\TRule{\True}{\Box}[2]}
          [\pFmla{\False}{\Box\lnot\formula{A}}{1},
            just ={\TRule{\True}{\lnot}[1]}
            [\pFmla{\False}{\lnot\formula{A}}{1.1},
              just = {\TRule{\False}{\Box}[4]}
              [\pFmla{\True}{\formula{A}}{1.1},
                  just ={\TRule{\False}{\lnot}[5]}, close]
            ]
          ]
        ]
      ]
    ]
  \end{oltableau}
  }{
  \begin{oltableau}
    [\pFmla{\True}{\Diamond \formula{A}}{1}, just = \TAss
      [\pFmla{\False}{\lnot\Diamond\lnot\formula{A}}{1}, just = \TAss
        [\pFmla{\True}{\formula{A}}{1.1}, just = {\TRule{\True}{\Diamond}[1]}
          [\pFmla{\True}{\Diamond\lnot\formula{A}}{1},
            just ={\TRule{\False}{\lnot}[2]}
            [\pFmla{\True}{\lnot\formula{A}}{1.1},
              just = {\TRule{\True}{\Diamond}[4]}
              [\pFmla{\False}{\formula{A}}{1.1},
                  just ={\TRule{\True}{\lnot}[5]}, close]
            ]
          ]
        ]
      ]
    ]
  \end{oltableau}
}}{
  \begin{oltableau}
    [\pFmla{\True}{\Diamond \formula{A}}{1}, just = \TAss,name=one
      [\pFmla{\False}{\Box \formula{A}}{1}, just = \TAss, name=two
        [\pFmla{\True}{\formula{A}}{1.1}, just = {\TRule{\True}{\Diamond}[1]}
          [\pFmla{\False}{\formula{A}}{1.1},
            just = {\TRule{\False}{\Box}[2]}, close]
        ]
      ]
    ]
  \end{oltableau}}

\end{document}
